% Generated draft: requires mathematical and source review.
\documentclass[11pt]{article}
\usepackage[margin=25mm]{geometry}
\usepackage{fontspec}
\setmainfont{Latin Modern Roman}
\usepackage{amsmath,amssymb,mathtools,unicode-math}
\setmathfont{Latin Modern Math}
\usepackage{enumitem,needspace,xurl,hyperref}
\setlength{\parindent}{0pt}
\setlength{\parskip}{5pt}
\setlength{\emergencystretch}{3em}
\newcommand{\N}{\mathbb{N}}
\newcommand{\Z}{\mathbb{Z}}
\newcommand{\Q}{\mathbb{Q}}
\newcommand{\R}{\mathbb{R}}
\newcommand{\C}{\mathbb{C}}
\newcommand{\F}{\mathbb{F}}
\newcommand{\A}{\mathbb{A}}
\newcommand{\PP}{\mathbb{P}}
\newcommand{\E}{\mathbb{E}}
\newcommand{\eps}{\varepsilon}
\newcommand{\dd}{\,\mathrm{d}}
\begin{document}
\section*{UnsolvedMath ID: 10600007}
\textbf{Source problem code:} catalogue entry
\subsection{Definitions and mathematical statement}
A long virtual knot diagram is an oriented immersed copy of \(\R\) in the plane, equal to a fixed straight line outside a compact set, with each double point designated either a classical over/under crossing or a virtual crossing. Its equivalence uses the intended virtual Reidemeister and detour moves relative to the fixed ends. A long flat virtual diagram replaces classical crossings by flat crossings and uses the corresponding flat virtual moves.

For a long flat diagram \(F\), traverse it from the incoming end. At each nonvirtual crossing assign the first encountered branch to pass over the second; keep virtual crossings virtual. Denote the resulting descending long virtual diagram by \(D(F)\). The source attributes to Turaev, in conversation, the assertion that \([D(F)]\) depends only on the long flat type of \(F\). Consequently, any appropriate virtual-knot invariant \(I\) defined on these long diagrams yields a long flat invariant \(F\mapsto I(D(F))\). If an invariant is defined only for closed knots, an appropriate closing convention would also have to be supplied.

The compound problem asks to:
\begin{enumerate}
\item enlarge the long-knot invariant structure of an unspecified cited paper, and determine whether that approach produces new classical-knot invariants;
\item combine the unspecified cited ideas with the cited biquandle construction to obtain stronger long-knot invariants;
\item calculate examples of the invariants \(I(D(F))\), and ultimately give a full classification of long flat knots, meaning complete invariants with an explicit description of their realizable values.
\end{enumerate}
The source also reports quaternion-invariant applications to long virtual knots by Bartholomew, Fenn, S. Kamada and N. Kamada, and points to another section for further questions. The actual constructions and references are missing. No particular meaning of stronger, new, or full invariant structure is specified, so these remain construction and classification tasks rather than fixed comparison inequalities.
\subsection{Short English statement}
Develop stronger long-knot invariants and possible classical-knot applications, combine the intended biquandle methods, compute invariants via Turaev's descending construction, and classify long flat virtual knots.
\subsection{Sources}
\begin{itemize}
\item \emph{Unsolved Problems in Virtual Knot Theory and Combinatorial Knot Theory}, revised survey, problem-list item 7, ``Long Knots and Long Flat Knots.'' \url{https://arxiv.org/abs/1409.2823}. Recorded access 2026-07-29; extraction source-tex. The statement credits Turaev's conversation and the quaternion work of Bartholomew, Fenn, S. Kamada and N. Kamada; several citation and section references are absent.
\item The supplied research summary lists incomplete references: \emph{Prime Decomposition and Non-Commutativity in the Monoid of Long Virtual Knots} (2013, arXiv:1311.xxxx); \emph{The intersection polynomials of a long virtual knot I/II} (2025, arXiv:2512.xxxx); \emph{Biquandle longitude invariant of long virtual knots} (2007, arXiv:0709.xxxx); \emph{Finite-Type Invariants of order one for long virtual knots} (2016, arXiv:1602.xxxx); and \emph{Band-Passes and Long Virtual Knot Concordance} (2016, arXiv:1603.xxxx). The printed identifiers are placeholders, not complete arXiv references.
\item Local UnsolvedMath registry, \texttt{data/unsolvedmath/problems.json}, permanent ID \(10600007\), exact source code \texttt{catalogue entry}; supplied partially\_solved/PARTIAL-PROGRESS status.
\end{itemize}
\end{document}
